\subsection{Q6: Attention}

% Vấn đề giải quyết
\subsubsection{Vấn đề giải quyết}
Bài toán yêu cầu cài đặt một khối \emph{self-attention một đầu} có mặt nạ nhân
quả (causal mask), là thành phần cốt lõi của kiến trúc Transformer và là bước
chuẩn bị cho mô hình sinh ký tự ở Q7. Đầu vào là một batch gồm $B$ chuỗi, mỗi
chuỗi có $T$ vị trí, mỗi vị trí được biểu diễn bởi một vector $C$ chiều; tức là
tensor $X \in \mathbb{R}^{B \times T \times C}$. Với mỗi vị trí $i$, khối
attention tính một vector đầu ra là tổng có trọng số của thông tin tại các vị
trí \emph{không đứng sau} $i$, nên đầu ra có cùng hình dạng
$O \in \mathbb{R}^{B \times T \times C}$.

Ràng buộc nhân quả là điều kiện bắt buộc: vị trí $i$ chỉ được nhìn các vị trí
$j \le i$. Nếu không, khi huấn luyện mô hình dự đoán ký tự tiếp theo, mô hình có
thể ``nhìn trộm'' chính ký tự cần dự đoán và không học được gì. Khác với Q4
(RNN đọc tuần tự từng ký tự), attention xử lý tất cả vị trí cùng lúc bằng phép
nhân ma trận, và mọi cặp vị trí đều tương tác trực tiếp thay vì thông qua một
trạng thái ẩn duy nhất.

% Kiến thức sử dụng
\subsubsection{Kiến thức sử dụng}

\textbf{Query, Key, Value.} Từ cùng một đầu vào $X$, ba tầng tuyến tính khác nhau
tạo ra ba tensor có cùng hình dạng:
\begin{equation}
Q = X W_Q + b_Q, \qquad K = X W_K + b_K, \qquad V = X W_V + b_V,
\end{equation}
trong đó $W_Q, W_K, W_V \in \mathbb{R}^{C \times C}$. Về trực giác, $Q_i$ là điều
vị trí $i$ ``đang tìm'', $K_j$ là điều vị trí $j$ ``có thể cung cấp'' để so
khớp, và $V_j$ là thông tin thực sự được lấy từ $j$ khi được chọn. Vì $Q, K, V$
đều sinh ra từ cùng một $X$ nên đây là \emph{self}-attention.

\textbf{Điểm tương đồng và hệ số chuẩn hoá.} Điểm giữa vị trí $i$ và $j$ là tích
vô hướng $Q_i \cdot K_j$. Với một batch, toàn bộ điểm được tính bằng một phép
nhân ma trận:
\begin{equation}
S = \frac{Q K^{\top}}{\sqrt{d_k}} \in \mathbb{R}^{B \times T \times T}, \qquad d_k = C.
\end{equation}
Phép chia cho $\sqrt{d_k}$ có lý do thống kê: nếu các thành phần của $Q_i$ và
$K_j$ độc lập, kỳ vọng 0 và phương sai 1, thì $Q_i \cdot K_j$ có phương sai $d_k$.
Khi $d_k$ lớn, điểm có độ lớn lớn khiến softmax bão hoà (gần one-hot) và gradient
rất nhỏ; chia cho $\sqrt{d_k}$ đưa phương sai về xấp xỉ 1. Với tensor 3 chiều,
phép chuyển vị $K^{\top}$ chỉ hoán đổi hai chiều cuối, giữ nguyên chiều batch
(trong cài đặt là \texttt{movedim(K, 1, 2)}).

\textbf{Mặt nạ nhân quả.} Mặt nạ là ma trận tam giác dưới
$\mathrm{tril}(\mathbf{1}_{T \times T})$, với phần tử $(i, j)$ bằng 1 khi $j \le i$.
Ví dụ với $T = 3$:
\begin{equation}
\mathrm{mask} = \begin{bmatrix} 1 & 0 & 0 \\ 1 & 1 & 0 \\ 1 & 1 & 1 \end{bmatrix}.
\end{equation}
Các điểm ở vị trí mặt nạ bằng 0 (tương lai) được gán $-\infty$ \emph{trước}
softmax:
\begin{equation}
\tilde{S}_{ij} = \begin{cases} S_{ij} & j \le i, \\ -\infty & j > i. \end{cases}
\end{equation}
Việc đặt mặt nạ trước softmax là quan trọng. Do $e^{-\infty} = 0$, các vị trí
tương lai có trọng số đúng bằng 0, mà mỗi hàng vẫn có tổng bằng 1. Nếu nhân mặt
nạ \emph{sau} softmax, các hàng sẽ không còn tổng bằng 1; còn nếu gán điểm bằng
0 thay vì $-\infty$, các vị trí tương lai vẫn có trọng số dương. Mỗi hàng luôn
còn ít nhất phần tử đường chéo ($j = i$) hữu hạn nên softmax không bao giờ sinh
\texttt{NaN}.

\textbf{Softmax theo hàng và tổng có trọng số.} Softmax được áp dụng trên chiều
cuối, tức là chiều của các Key, để mỗi vị trí $i$ có một phân phối trọng số trên
các vị trí được phép nhìn:
\begin{equation}
A_{ij} = \frac{e^{\tilde{S}_{ij}}}{\sum_{j'} e^{\tilde{S}_{ij'}}}, \qquad
O = A\,V.
\end{equation}
Để thấy hình dạng của $A$, xét trường hợp đặc biệt mọi điểm hữu hạn bằng nhau
(với $T = 3$):
\begin{equation}
A = \begin{bmatrix} 1 & 0 & 0 \\ \tfrac12 & \tfrac12 & 0 \\ \tfrac13 & \tfrac13 & \tfrac13 \end{bmatrix}.
\end{equation}
Hàng đầu tiên luôn là $[1, 0, 0]$ vì vị trí đầu chỉ được nhìn chính nó. Ma trận
$A$ luôn là tam giác dưới với mỗi hàng có tổng bằng 1, và đây là hai tính chất
dễ dùng để kiểm tra nhanh một cài đặt.

\textbf{Mặt nạ là buffer, không phải tham số.} Trong \texttt{\_\_init\_\_}, mặt nạ
được đăng ký bằng \texttt{register\_buffer}. Buffer đi theo mô hình khi chuyển
thiết bị và được lưu trong \texttt{state\_dict}, nhưng không có gradient và không
bị optimizer cập nhật. Số tham số học được của khối chỉ gồm ba tầng tuyến tính:
$3(C^2 + C)$.

\textbf{Chi phí.} Ma trận điểm có $T \times T$ phần tử mỗi chuỗi, nên chi phí tính
toán và bộ nhớ tăng theo bình phương độ dài chuỗi ($O(T^2 C)$ phép tính). Đổi lại
mọi cặp vị trí tương tác trực tiếp và toàn bộ vị trí được tính song song.

% Implement, công thức
\subsubsection{Cài đặt}

\textbf{Khởi tạo.} Lớp \texttt{Attention(layer\_size, block\_size)} trong
\texttt{models.py} đã định nghĩa sẵn ba tầng \texttt{Linear(layer\_size,
layer\_size)} tên \texttt{q\_layer}, \texttt{k\_layer}, \texttt{v\_layer}, và mặt
nạ $\mathrm{tril}(\mathbf{1})$ kích thước \texttt{block\_size} $\times$
\texttt{block\_size} (đổi hình dạng thành $1 \times 1 \times \texttt{block\_size}
\times \texttt{block\_size}$). Ở đây \texttt{layer\_size} chính là $C$ và cũng
là $d_k$ (mô hình dùng một đầu, không chia nhiều đầu). Autograder yêu cầu mỗi
tầng phải dùng đúng cho ma trận tương ứng của nó.

\textbf{Hàm \texttt{forward(input)}} gồm các bước, kèm hình dạng tensor ở mỗi
bước ($B$: batch, $T$: độ dài chuỗi thực tế, $C$: \texttt{layer\_size}):
\begin{table}[htbp]
\caption{Các bước của \texttt{forward} và hình dạng tensor}
\label{tab:q6_shapes}
\centering
\begin{tabular}{clc}
\toprule
Bước & Phép tính & Hình dạng \\
\midrule
1 & $Q$, $K$, $V$ từ \texttt{q\_layer}, \texttt{k\_layer}, \texttt{v\_layer} & $(B, T, C)$ \\
2 & $S = \mathrm{matmul}(Q, \mathrm{movedim}(K, 1, 2)) / \sqrt{C}$ & $(B, T, T)$ \\
3 & Gán $-\infty$ tại các vị trí mặt nạ bằng 0 (\texttt{masked\_fill}) & $(B, T, T)$ \\
4 & \texttt{softmax} theo chiều cuối & $(B, T, T)$ \\
5 & $O = \mathrm{matmul}(A, V)$ & $(B, T, C)$ \\
\bottomrule
\end{tabular}
\end{table}

\textbf{Lưu ý về chiều khi áp mặt nạ.} Mặt nạ cắt theo độ dài thực tế
(\texttt{mask[:, :, :T, :T]}) có hình dạng $(1, 1, T, T)$. Khi \texttt{masked\_fill}
broadcast mặt nạ này với $S$ hình dạng $(B, T, T)$, kết quả có thêm một chiều đầu
tiên và thành $(1, B, T, T)$. Vì vậy cần lấy phần tử \texttt{[0]} để quay về
$(B, T, T)$ trước khi qua softmax và nhân với $V$. Đây là lý do dòng gợi ý trong
đề bài kết thúc bằng \texttt{[0]}.

\textbf{Kiểm tra.} Autograder (\texttt{check\_attention}) so sánh đầu ra của khối
với đáp án mong đợi bằng \texttt{torch.isclose}; cài đặt trên đã vượt qua kiểm
tra này với số điểm 2/2. Q6 không có ngưỡng độ chính xác hay quá trình huấn
luyện như Q3 và Q4, vì đây là một khối tính toán xác định được kiểm tra trực
tiếp bằng đầu ra. Hai tính chất ở trên (ma trận $A$ là tam giác dưới, mỗi hàng
tổng bằng 1) là cách kiểm tra bổ sung khi gỡ lỗi.